\section{Hidden Response and Interaction Curvature}
\label{sec:interaction}

The algebraic reduction becomes a differential theory when the canonical
minimizer varies smoothly.  We fix a local product chart throughout this
section, and every ordinary second derivative below is taken in that chart.
The Schur formula is therefore a coordinate Hessian identity.  It defines an
intrinsic bilinear form at a full reduced critical point
\(D_z\bar E=0\), or after ordinary derivatives are replaced by covariant
derivatives for specified connections.  Fiber stationarity
\(E_\sigma=0\) alone does not make the visible coordinate Hessian a tensor.
The mechanism amplitudes below carry a fixed affine structure, under which
their Hessian transforms by congruence.

\subsection{Smooth response and effective curvature}

Let \(z=(y,u)\) range over a finite-dimensional manifold locally represented
in Euclidean coordinates, and let \(\Sigspace\) be a finite-dimensional
smooth manifold with local hidden coordinate \(\sigma\).  Consider
\[
  E:(z,\sigma)\longmapsto E(z,\sigma)\in\mathbb R.
\]

\begin{assumption}[Smooth stable reduction]
\label{ass:smooth-reduction}
For some \(r\ge2\), the energy is \(C^{r+1}\).  It is locally uniformly
inf-compact in \(\sigma\): for every compact parameter set \(K\) and every
\(c\in\mathbb R\), the union of the fiber sublevel sets
\[
  \{\sigma:\exists z\in K,\ E(z,\sigma)\le c\}
\]
is relatively compact.  Every \(z\) has a unique interior minimizer
\(\sigma_\star(z)\), and
\[
  H(z):=E_{\sigma\sigma}(z,\sigma_\star(z))
\]
is positive definite.
\end{assumption}

\begin{theorem}[Hidden response and Schur reduced Hessian]
\label{thm:response-schur}
Under Assumption~\ref{ass:smooth-reduction}, the minimizer section and reduced energy
\[
  \sigma_\star(z)=\argmin_\sigma E(z,\sigma),
  \qquad
  \bar E(z)=E(z,\sigma_\star(z))
\]
are \(C^r\).  At the canonical lift,
\begin{align}
  D_z\sigma_\star&=-H^{-1}E_{\sigma z},
  \label{eq:hidden-response}\\
  D_z\bar E&=E_z,
  \label{eq:envelope}\\
  D^2_{zz}\bar E
  &=E_{zz}-E_{z\sigma}H^{-1}E_{\sigma z}.
  \label{eq:schur-effective-hessian}
\end{align}
If \(\sigma=(\alpha,\beta)\), let \(\mathbb H\) denote the full Hessian in
\((z,\alpha,\beta)\) at the canonical lift.  When its joint hidden block is
positive definite, sequential and joint second-order elimination agree:
\begin{equation}
  \Schur_{(\alpha,\beta)}\mathbb H
  =\Schur_\alpha\bigl(\Schur_\beta\mathbb H\bigr).
  \label{eq:schur-associativity}
\end{equation}
\end{theorem}

The response formula is familiar from parametric optimization and variable
projection \cite{golub1973variableprojection,bonnans2000perturbation}.  Here
it supplies the local law inside a globally defined geometric action bundle
and composes with the optimizer hierarchy in \cref{thm:algebraic-reduction}.

\subsection{Affine perturbations before reduction}

Let \(u\in\Uspace\subset\mathbb R^p\) denote continuous mechanism amplitudes.
Suppose
\begin{equation}
  E(y,\sigma,u)
  =E_0(y,\sigma)+\sum_{i=1}^p u_i\Delta_i(y,\sigma).
  \label{eq:affine-interventions}
\end{equation}
At the canonical lift define
\begin{equation}
  G:T_u\Uspace\to T_\sigma^*\Sigspace,
  \qquad
  Ga=\sum_{i=1}^p a_iD_\sigma\Delta_i.
  \label{eq:interaction-map}
\end{equation}

\begin{theorem}[Reduction-induced interaction curvature]
\label{thm:interaction-curvature}
Under Assumption~\ref{ass:smooth-reduction} and
\cref{eq:affine-interventions},
\begin{equation}
  \boxed{
  D^2_{uu}\bar E=-G^*H^{-1}G\preceq0.}
  \label{eq:interaction-curvature}
\end{equation}
Equivalently,
\begin{equation}
  \partial^2_{u_i u_j}\bar E
  =-
  \left\langle
  D_\sigma\Delta_i,
  H^{-1}D_\sigma\Delta_j
  \right\rangle.
  \label{eq:pair-interaction-curvature}
\end{equation}
Thus the self-curvature of every mechanism is nonpositive.  A pointwise mixed
interaction vanishes exactly when the corresponding vertical perturbations
are orthogonal in the inverse-Hessian metric.  Mixed entries may have either
sign.
\end{theorem}

Even without smoothness, \(u\mapsto\bar E(y,u)\) is concave wherever finite,
because it is a pointwise infimum of affine functions.  Smoothness identifies
the exact curvature responsible for that concavity.

The affine hypothesis is essential for the sign.  With general smooth
\(u\)-dependence, \cref{eq:schur-effective-hessian} gives
\begin{equation}
  D^2_{uu}\bar E
  =E_{uu}-E_{u\sigma}H^{-1}E_{\sigma u},
  \label{eq:nonaffine-interaction-curvature}
\end{equation}
so the direct curvature \(E_{uu}\) need not be negative semidefinite.  The
quadratic form is covariant under an affine reparameterization
\(u=T\tilde u+u_0\), transforming as
\(D^2_{\tilde u\tilde u}\bar E=T^*D^2_{uu}\bar E\,T\).  Individual mixed
entries therefore have meaning only after the mechanism basis, direction,
and units are fixed; all comparisons in this paper refer to the specified
amplitude protocol.

Factorial finite differences are iterated integrals of the corresponding
mixed derivatives; \cref{cor:factorial-contrast} states the general identity.
In particular,
\begin{equation}
  \delta_{\{i,j\}}\bar E
  =-
  \int_0^1\!\int_0^1
  \left\langle D_\sigma\Delta_i,
  H^{-1}D_\sigma\Delta_j\right\rangle
  \,\dd u_i\,\dd u_j,
  \label{eq:pair-factorial-integral}
\end{equation}
where all quantities are evaluated along the canonical minimizer with the
other amplitudes fixed at zero.
The integral may vanish by cancellation even when pointwise interaction does
not.  Reduction must precede finite differencing; the counterexample and the
exact commuting special case are given in \cref{app:finite-contrasts}.

\subsection{Action-fiber specialization}

The global action chart makes the same response tensor computational.  For
fixed \(B\in\Bact\), perturb the canonical fiber objective by
\begin{equation}
  F_B(\Sigma,u)
  =f_B(\Sigma)+\sum_{i=1}^p u_i\Delta_i(B,\Sigma),
  \qquad \Sigma\in\mathscr P_{d-m}.
  \label{eq:perturbed-action-fiber}
\end{equation}
Let \(H_B=\operatorname{Hess}f_B(\Sigma_\star)\) and
\(G_Ba=\sum_i a_iD_\Sigma\Delta_i(B,\Sigma_\star)\).  Strong convexity of
the action objective gives \(H_B\succeq\operatorname{Id}\) as a self-adjoint
operator with respect to the AIRM metric.  For an intrinsic perturbation,
we require the natural complement-frame covariance
\(\Delta_{i,VR}(B,R^\top\Sigma R)=\Delta_{i,V}(B,\Sigma)\); otherwise the
chosen \(\Delta_i\) deliberately selects a hidden frame.
Derivatives with respect to \(B\) use the natural affine structure of the open
matrix domain \(\Bact\subset\mathbb R^{d\times m}\).

\begin{corollary}[Interaction response of the canonical action controller]
\label{cor:action-interaction-response}
Assume the \(\Delta_i\) are \(C^2\) near \(\Sigma_\star\).  For sufficiently
small \(u\), there is a unique critical point \(\Sigma_\star(B,u)\) near
\(\Sigma_\star\); it is a strict local minimizer.  Define the corresponding
local value branch by
\(\bar F_B^{\rm loc}(u)=F_B(\Sigma_\star(B,u),u)\).  Then
\begin{equation}
  D_u\Sigma_\star(B,0)=-H_B^{-1}G_B,
  \qquad
  D^2_{uu}\bar F_B^{\rm loc}(0)
  =-G_B^*H_B^{-1}G_B.
  \label{eq:action-interaction-response}
\end{equation}
The same inverse vertical Hessian \(H_B^{-1}\) governs the hidden-coordinate
response to the visible action \(B\).  The full section derivative also
contains the chart's direct term:
\[
  D_Bs_A=D_B\Phi_A+D_\Sigma\Phi_A[D_B\Sigma_\star].
\]
Thus interaction curvature, action susceptibility, and the residual solver
are local response, base response, and computation for one canonical fiber
problem.
\end{corollary}

This response statement is local in the perturbation.  A global perturbed
minimum requires separate coercivity or growth control on the \(\Delta_i\).

\begin{proof}
At \((\Sigma_\star,0)\), the derivative in \(\Sigma\) of the stationarity
equation is the invertible operator \(H_B\).  The implicit-function theorem
therefore gives the unique nearby critical branch; continuity of its Hessian
makes it a strict local minimum after shrinking the \(u\)-neighborhood.
Differentiating stationarity and then the local value branch gives
\cref{eq:action-interaction-response}.  The covariance law makes this branch
independent of complement coordinates.  Differentiation with respect to \(B\)
uses the same stationarity equation with \(E_{\Sigma B}\) in place of \(G_B\).
\end{proof}
